%% parts/p06/descriptive-set-theory.tex --- Phase-4 P6 depth fill
%%
%% Register coverage: D-P7-013, N-a-046/047, A-068, A-135.
%% Source accounting: v1.1 Part VII has no descriptive-set-theory development
%% to migrate.  Every result here is new, register-driven, and
%% citation-backed; no source proof status is promoted.

\chapter{Descriptive Set Theory}
\label{sec:p06-descriptive-set-theory}
\chaptermark{Descriptive Set Theory}

This chapter classifies definable subsets of well-behaved topological spaces
and calibrates their regularity properties.  It is the classical, boldface
counterpart of the effective hierarchy treated in Part~4; the following
chapter supplies the separate game-theoretic calibration.
The standing modern references are \cite{Kechris1995,Moschovakis2009}; the
historical attribution for the analytic--coanalytic characterization is
\cite{Suslin1917}.

\section{Polish spaces and the Borel hierarchy}
\label{sec:p06-polish-borel}

\begin{hbdefinition}{def:p06-polish-space}{Polish space and Baire space}
A topological space \(X\) is \emph{Polish} iff it has a countable dense subset
and admits a complete metric whose open balls generate the given topology.
Thus the adjective names a property of a topological space, not a preferred
metric: two different compatible complete metrics may witness the same Polish
topology.  Completeness alone would omit separability, while separability alone
would omit the complete compatible metric; both clauses are needed below for
the representation and regularity theorems.

The set of all functions from \(\omega\) to \(\omega\), equipped with the
product topology obtained from the discrete topology on \(\omega\), is
\emph{Baire space}.  In this chapter the expression \(\omega^\omega\) names
that function space, not ordinal exponentiation.  Baire space is the witness
space used in projections defining analytic and projective sets.
\index{Polish space}
\index{Baire space}
\end{hbdefinition}

\begin{hbdefinition}{def:p06-borel-hierarchy}{Boldface Borel hierarchy}
A \emph{pointclass} on \(X\) is a metalinguistically specified family of
subsets of \(X\).  This term is needed because the hierarchy records closure
and definability uniformly as the ambient Polish space varies; a pointclass is
not itself a set inside \(X\).

Let \(X\) be Polish.  For every nonzero countable ordinal
\(\alpha<\omega_1\), define
\[
\begin{aligned}
  \boldsymbol{\Sigma}^{0}_{1}(X)
    &\defeq \{O\subseteq X:O\text{ is open}\},\\
  \boldsymbol{\Pi}^{0}_{\alpha}(X)
    &\defeq \{X\setminus A:A\in
       \boldsymbol{\Sigma}^{0}_{\alpha}(X)\},\\
  \boldsymbol{\Sigma}^{0}_{\alpha}(X)
    &\defeq
      \left\{\bigcup_{m\in\omega}A_m:
      (\forall m)\,
      A_m\in
      \bigcup_{1\leq\beta<\alpha}
      \boldsymbol{\Pi}^{0}_{\beta}(X)\right\}
      \quad(1<\alpha<\omega_1),\\
  \boldsymbol{\Delta}^{0}_{\alpha}(X)
    &\defeq
      \boldsymbol{\Sigma}^{0}_{\alpha}(X)
      \cap\boldsymbol{\Pi}^{0}_{\alpha}(X).
\end{aligned}
\]
The \emph{Borel sets} are
\[
  \operatorname{Bor}(X)
  \defeq \bigcup_{1\leq\alpha<\omega_1}
  \boldsymbol{\Sigma}^{0}_{\alpha}(X),
\]
equivalently the least sigma-algebra containing the open subsets of \(X\).
The superscript \(0\) distinguishes the Borel hierarchy from the projective
hierarchy below, and boldface records that no effective restriction is placed
on the codes.  The level index is a countable ordinal, not the number of
quantifier blocks in the lightface arithmetical hierarchy
\cite[Sections~11--13]{Kechris1995}.
\index{pointclass}
\index{Borel hierarchy}
\index{Borel set}
\index{Sigma-zero-alpha@\(\boldsymbol{\Sigma}^{0}_{\alpha}\)}
\index{Pi-zero-alpha@\(\boldsymbol{\Pi}^{0}_{\alpha}\)}
\index{Delta-zero-alpha@\(\boldsymbol{\Delta}^{0}_{\alpha}\)}
\end{hbdefinition}

\begin{hbtheorem}{thm:p06-borel-exhaustion}{cited=Kechris1995}{Borel-rank exhaustion}
For every Polish space \(X\), each Borel subset of \(X\) belongs to
\(\boldsymbol{\Sigma}^{0}_{\alpha}(X)\) and to
\(\boldsymbol{\Pi}^{0}_{\beta}(X)\) for some countable ordinals
\(\alpha,\beta<\omega_1\).  At every fixed level,
\(\boldsymbol{\Sigma}^{0}_{\alpha}\) is closed under countable unions,
\(\boldsymbol{\Pi}^{0}_{\alpha}\) is closed under countable intersections,
and complementation interchanges the paired classes.  The theorem does not
put one uniform countable bound on the ranks of all Borel subsets of an
uncountable Polish space.
\end{hbtheorem}

\section{Analytic and projective pointclasses}
\label{sec:p06-projective-pointclasses}

\begin{hbdefinition}{def:p06-projective-pointclasses}{Analytic, coanalytic, and projective sets}
Let \(X\) be Polish.  The first-coordinate \emph{projection} of
\(C\subseteq X\times\omega^\omega\) is
\(\{x\in X:(\exists y\in\omega^\omega)\,(x,y)\in C\}\).  A set
\(A\subseteq X\) is \emph{analytic}, written
\(A\in\boldsymbol{\Sigma}^{1}_{1}(X)\), iff it is the projection of some
closed \(C\subseteq X\times\omega^\omega\).  A set is \emph{coanalytic}, written
\(A\in\boldsymbol{\Pi}^{1}_{1}(X)\), iff its complement is analytic, and
\[
  \boldsymbol{\Delta}^{1}_{1}(X)
  \defeq
  \boldsymbol{\Sigma}^{1}_{1}(X)
  \cap
  \boldsymbol{\Pi}^{1}_{1}(X).
\]
The empty set is covered by the projection definition; when \(A\ne\varnothing\),
analyticity is equivalently representability as a continuous image of Baire
space.

For \(n\geq1\), the higher \emph{projective pointclasses} are defined by
\[
\begin{aligned}
  \boldsymbol{\Sigma}^{1}_{n+1}(X)
    &\defeq
      \left\{A\subseteq X:
        \begin{array}{l}
        (\exists B\in\boldsymbol{\Pi}^{1}_{n}
          (X\times\omega^\omega))\\[-2pt]
        A=\{x\in X:(\exists y\in\omega^\omega)\,(x,y)\in B\}
        \end{array}\right\},\\
  \boldsymbol{\Pi}^{1}_{n+1}(X)
    &\defeq
      \{X\setminus A:A\in
        \boldsymbol{\Sigma}^{1}_{n+1}(X)\},\\
  \boldsymbol{\Delta}^{1}_{n}(X)
    &\defeq
      \boldsymbol{\Sigma}^{1}_{n}(X)
      \cap\boldsymbol{\Pi}^{1}_{n}(X).
\end{aligned}
\]
These are families of subsets of Polish spaces generated by projection and
complement.  They are not the formula classes of the L\'evy hierarchy, and
they are not the lightface analytical classes of Part~4.  Their additional
closure under real projection is the distinction needed for Suslin's theorem,
the regularity boundary, and the projective game axiom defined in the
following chapter
\cite[Sections~25, 37]{Kechris1995}.
\index{projection!descriptive set theory}
\index{analytic set}
\index{coanalytic set}
\index{projective hierarchy!boldface}
\index{Sigma-one-n@\(\boldsymbol{\Sigma}^{1}_{n}\)}
\index{Pi-one-n@\(\boldsymbol{\Pi}^{1}_{n}\)}
\index{Delta-one-n@\(\boldsymbol{\Delta}^{1}_{n}\)}
\end{hbdefinition}

\begin{hbtheorem}{thm:p06-suslin}{cited=Kechris1995}{Suslin's theorem}
Let \(X\) be Polish and \(A\subseteq X\).  Then
\[
  A\in\operatorname{Bor}(X)
  \quad\Leftrightarrow\quad
  A\in
  \boldsymbol{\Sigma}^{1}_{1}(X)
  \cap
  \boldsymbol{\Pi}^{1}_{1}(X).
\]
Equivalently,
\(\operatorname{Bor}(X)=\boldsymbol{\Delta}^{1}_{1}(X)\).
The historical source is \cite{Suslin1917}; the status citation supports the
modern Polish-space formulation.
\index{Suslin theorem@Suslin's theorem}
\end{hbtheorem}

\begin{hbdefinition}{def:p06-regularity-properties}{Regularity properties}
Let \(A\subseteq X\), where \(X\) is Polish.
\begin{enumerate}
  \item A set is \emph{meager} iff it is a countable union of nowhere-dense
        sets.  The set \(A\) has the \emph{Baire property} iff
        \((A\setminus O)\cup(O\setminus A)\) is meager for some open
        \(O\subseteq X\).
  \item A subset \(P\subseteq X\) is \emph{perfect} iff it is nonempty,
        closed, and has no isolated point.  The set \(A\) has the
        \emph{perfect-set property} iff it is countable or contains a perfect
        subset.
  \item The set \(A\) is \emph{universally measurable} iff it belongs to the
        completion of the Borel sigma-algebra for every sigma-finite Borel
        measure on \(X\).
\end{enumerate}
These are three logically distinct properties of subsets of a Polish space;
the phrase \emph{regularity properties} below is only a metalinguistic
abbreviation for this displayed list.  None is part of the definition of a
projective pointclass.  The distinction is required because ZFC proves
different rows of the list at different pointclass levels
\cite[Sections~21, 29, 35]{Kechris1995}.
\index{Baire property}
\index{meager set}
\index{perfect-set property}
\index{universal measurability}
\end{hbdefinition}

\begin{hbtheorem}{thm:p06-analytic-regularity}{cited=Kechris1995}{ZFC regularity of analytic sets}
Let \(X\) be Polish.  Every analytic subset of \(X\) has the Baire property,
the perfect-set property, and is universally measurable.  Consequently, an
uncountable analytic set contains a nonempty perfect subset and has cardinality
\(2^{\aleph_0}\).  Every Borel set has the same three properties because every
Borel set is analytic.
\end{hbtheorem}

\begin{hbtheorem}{thm:p06-coanalytic-boundary}{cited=Kechris1995}{The coanalytic regularity boundary}
Every coanalytic subset of a Polish space has the Baire property and is
universally measurable.  Moreover, \(\mathsf{ZFC}+V=L\) proves that there is
an uncountable coanalytic set of reals with no perfect subset.  Since
\(\mathsf{ZFC}+V=L\) is equiconsistent with \(\mathsf{ZFC}\)
\parencite{Godel1938}, it follows, assuming
\(\mathsf{Con}(\mathsf{ZFC})\), that \(\mathsf{ZFC}\) does not prove that every
uncountable coanalytic set has a perfect subset.  Thus complementation transfers
the first and third regularity properties from analytic sets, but not, over
\(\mathsf{ZFC}\) alone, the perfect-set property.
\end{hbtheorem}

\begin{proofboundary}{def:p06-projective-regularity-boundary}{Projective regularity boundary}
The preceding ZFC theorems do not extend to a blanket assertion that every
projective set has the three regularity properties.  At coanalytic level the
perfect-set boundary is already visible, and at higher projective levels the
table below records no unconditional ZFC check.  Additional game axioms,
defined in the following chapter, can supply stronger regularity, but they
belong in the extra-hypothesis column and must not be read back into ZFC.
\end{proofboundary}

The word \emph{lightface} denotes the effective, parameter-free versions of
these pointclasses, obtained from recursive codes rather than arbitrary real
parameters.  That adjective therefore marks a metalinguistic restriction on
codes, not a new topological operation.  The lightface analytical and
hyperarithmetical hierarchies are treated in Part~4; the present sentence is
a cross-part bridge, and no separate effective descriptive-set-theory
development is added here
\cite{Sacks1990,Moschovakis2009}.
\index{projective hierarchy!lightface pointer}

\begin{figure}[!htbp]
\centering
\begin{tikzpicture}[
  x=1cm,y=1cm,
  pc/.style={draw,rounded corners,fill=blue!5,minimum width=2.35cm,
             minimum height=.70cm,align=center,font=\scriptsize},
  borel/.style={draw,rounded corners,fill=green!7,minimum width=2.25cm,
                minimum height=.70cm,align=center,font=\scriptsize},
  op/.style={-{Latex[length=2.1mm]},thick},
  unionop/.style={-{Latex[length=2.1mm]},thick,densely dotted},
  iterate/.style={-{Latex[length=2.1mm]},thick,dashed},
  comp/.style={<->,>=Latex,thick},
  note/.style={font=\scriptsize,align=center}
]
  \node[borel] (s01) at (-2.7,5.8)
    {\(\boldsymbol{\Sigma}^{0}_{1}\)\\open};
  \node[borel] (p01) at (-2.7,4.55)
    {\(\boldsymbol{\Pi}^{0}_{1}\)\\closed};
  \node[borel] (sa) at (2.7,5.8)
    {\(\boldsymbol{\Sigma}^{0}_{\alpha}\)\\countable unions};
  \node[borel] (pa) at (2.7,4.55)
    {\(\boldsymbol{\Pi}^{0}_{\alpha}\)\\complements};
  \draw[comp] (s01) -- (p01);
  \draw[comp] (sa) -- (pa);
  \draw[iterate] (s01) -- (sa);
  \draw[unionop] (p01.east) -- (sa.south west);

  \node[borel,minimum width=5.2cm] (bor) at (0,3.25)
    {\(\operatorname{Bor}(X)=\boldsymbol{\Delta}^{1}_{1}(X)\)\\[-1pt]
     \(\displaystyle
      =\bigcup_{1\leq\alpha<\omega_1}
       \boldsymbol{\Sigma}^{0}_{\alpha}(X)\)};
  \draw[op] (pa.south) -- (bor.north east);

  \node[pc] (s11) at (-2.0,1.75)
    {\(\boldsymbol{\Sigma}^{1}_{1}\)\\analytic};
  \node[pc] (p11) at (2.0,1.75)
    {\(\boldsymbol{\Pi}^{1}_{1}\)\\coanalytic};
  \draw[op] (bor.south west) -- (s11.north);
  \draw[op] (bor.south east) -- (p11.north);
  \draw[comp] (s11) -- (p11);

  \node[pc] (s12) at (-2.0,-.05)
    {\(\boldsymbol{\Sigma}^{1}_{2}\)};
  \node[pc] (p12) at (2.0,-.05)
    {\(\boldsymbol{\Pi}^{1}_{2}\)};
  \draw[op] (p11.south west) -- (s12.north east);
  \draw[comp] (s12) -- (p12);

  \node[pc] (s1n) at (-2.0,-1.85)
    {\(\boldsymbol{\Sigma}^{1}_{n+1}\)};
  \node[pc] (p1n) at (2.0,-1.85)
    {\(\boldsymbol{\Pi}^{1}_{n+1}\)};
  \draw[iterate] (p12.south west) -- (s1n.north east);
  \draw[comp] (s1n) -- (p1n);

  \node[note] at (0,-2.75)
    {Boldface pointclasses: arbitrary real parameters are allowed.\\
     Lightface effective pointclasses are only cross-referenced to Part~4.};
\end{tikzpicture}
\caption{Borel and projective hierarchy operations.
The upper row abbreviates the transfinite Borel progression through nonzero
countable \(\alpha\): double arrows mean complement, the dotted diagonal means
countable union from lower dual levels, and the dashed arrow means iteration.
Arrows from the Borel box into the first projective pair are inclusions;
downward diagonals are projections and horizontal double arrows are
complements.  No arrow asserts a collapse of projective sets to Borel sets.}
\label{fig:p06-borel-projective}
\end{figure}

\begin{landscape}
\begingroup\footnotesize
\setlength{\tabcolsep}{4pt}
\begin{longtable}{@{}p{.13\linewidth}p{.20\linewidth}p{.20\linewidth}
  p{.24\linewidth}p{.17\linewidth}@{}}
\caption{Borel and projective pointclasses.  The fourth
column contains only \(\mathsf{ZFC}\) consequences; game results and
large-cardinal assumptions are isolated in the final column.}
\label{tab:p06-borel-projective}\\
\toprule
Pointclass & Definition & Closure recorded here & ZFC regularity &
Game result or extra hypothesis \\
\midrule
\endfirsthead
\caption[]{Borel and projective pointclasses.  The fourth
column contains only \(\mathsf{ZFC}\) consequences; game results and
large-cardinal assumptions are isolated in the final column (continued)}\\
\toprule
Pointclass & Definition & Closure recorded here & ZFC regularity &
Game result or extra hypothesis \\
\midrule
\endhead
\midrule\multicolumn{5}{r}{\emph{continued on next page}}\\
\endfoot
\endlastfoot
Borel
\(\operatorname{Bor}(X)=\boldsymbol{\Delta}^{1}_{1}(X)\) &
Transfinite countable union and complement from open sets &
Sigma-algebra; continuous preimages &
Baire property, perfect-set property, universal measurability &
\(\mathsf{ZFC}\); exact game statement in
Theorem~\ref{thm:p06-borel-determinacy} \\
Analytic
\(\boldsymbol{\Sigma}^{1}_{1}(X)\) &
Projection of a closed subset of \(X\times\omega^\omega\) &
Continuous images, Borel preimages, countable unions and intersections &
Baire property, perfect-set property, universal measurability &
Additional hypothesis; exact calibration in
Definition~\ref{def:p06-ad-pd} \\
Coanalytic
\(\boldsymbol{\Pi}^{1}_{1}(X)\) &
Complement of an analytic set &
Complement dual; Borel preimages, countable unions and intersections &
Baire property and universal measurability; no unconditional
\(\mathsf{ZFC}\) perfect-set assertion &
Under \(V=L\) there is an uncountable coanalytic counterexample to the
perfect-set property; additional determinacy hypotheses are calibrated in
Definition~\ref{def:p06-ad-pd} \\
Higher projective
\(\boldsymbol{\Sigma}^{1}_{n+1},
 \boldsymbol{\Pi}^{1}_{n+1}\) &
Projection from \(\boldsymbol{\Pi}^{1}_{n}\), then complement &
Continuous preimages and countable unions/intersections; complement swaps
the paired classes &
No blanket ZFC assertion of all three regularity properties &
No ZFC check; sufficient large-cardinal hypothesis in
Theorem~\ref{thm:p06-martin-steel} \\
All sets of reals (full determinacy) &
Payoffs \(A\subseteq\omega^\omega\), with no pointclass restriction &
Not a pointclass closure claim &
No blanket \(\mathsf{ZFC}\) determinacy assertion &
The assertion that every length-\(\omega\) integer game is determined implies
that the reals cannot be well ordered and hence contradicts
\(\mathsf{AC}\); see Theorem~\ref{thm:p06-ad-versus-choice} \\
\bottomrule
\end{longtable}
\endgroup
\end{landscape}
\thispagestyle{empty}% suppress the intentionally blank verso before Chapter 40
